\noindent This work investigates the impact of in-season roster trades on the performance of \textit{National Basketball Association} (NBA) teams, treating a trade as a multidimensional event rather than something reducible to a single indicator. Four hypotheses were operationalized and tested. The first addresses the association between the variation in on-court salary inequality — measured by the Gini Index weighted by minutes played, selected among six competing salary inequality metrics — and a team's subsequent performance. The second examines the association between the functional fit of an acquired player, diagnosed through the Relative Contribution methodology derived from \textit{Win Probability Added} (WPA), and the trade's short-term impact. The third assesses the association between the experience of acquired players, measured by their average age at the time of the transaction, and that same impact. The fourth investigates the association between the severity of the positional gap filled and the size of the player package negotiated. Data were collected through a layered data architecture (\textit{Medallion Lakehouse}), drawn from public APIs and historical NBA sources, covering thirteen seasons, from 2012-13 to 2024-25. Hypotheses were tested using Spearman correlation and the Mann-Whitney \textit{U} test, on samples ranging from 282 to 370 cases depending on the test. Results show statistically significant, though weak, support for three of the four hypotheses. Salary inequality is negatively associated with subsequent performance ($\rho = -0.189$; $p = 0.0014$), moderately favoring Equity Theory over Tournament Theory and concentrated among offensively-oriented teams. The experience of acquired players is positively associated with the trade's short-term impact ($\rho = 0.110$; $p = 0.0346$). The size of the negotiated package, in turn, tracks the severity of the gap filled ($\rho = -0.280$; $p < 0.0001$). The functional fit of the acquired player, however, finds no statistical support in its original categorical form, although a weaker signal survives in its continuous version. Taken together, the four results support the reading that a roster trade is indeed a multidimensional event, in which no single dimension alone determines the success of the move.

\vspace{0.5cm}
\noindent\textbf{Keywords}: Roster trades. Salary inequality. Equity Theory. Sports performance. NBA.
